\documentclass{article}
\usepackage{/globals/laws}
\usepackage{/globals/de}
\usepackage{/globals/anki}
\ankiprefix{gesetze}
\ankitopic{Gesetze}
\ankishow
\begin{document}
\section{Gesetze}
Ein Gesetz ist eine thematisch sortierte Sammlung an unterschiedlichen Regelungen, Normen und Vorschriften. Gesetze haben einen Namen, die Gesetzesbezeichnung, eine Kurzbezeichnung und eine Abkürzung.

Sie sind in unterschiedliche Paragraphen, oder im Fall des Grundgesetzes oder des EU-Rechts in Artikel, aufgeteilt. In Gesetzestexten sind Paragraphen am Paragraphnesymbol \sqe{\para{1}} und Artikel an der Kurzform \sqe{\art{1}} zu erkennen.

Diese können in Absätze \sqe{\labsa{1}} unterteilt werden. Absätze bestehen oft aus mehreren Sätzen. Es können auch Nummern \sqe{\lnr{1}} und Buchstaben \sqe{\llit{a}} vorkommen sein.

\ankinote{aufteilung}{Aufteilung}{Gesetz, Paragraph/Artikel, Absatz, Satz, Nummber/Buchstabe}

\subsection{Zitieren}
Beim Zitieren wird der größe nach sortiert, von der grobsten Angabe zur genausten; außer der Name des Gesetzes selbst, dieses kommt immer als letztes. Somit ist die Reihenfolge: Paragraph/Artikel, Absatz, Satz, Nummber/Buchstabe, Gesetz.

Genutzt werden die Abkürzungen und Symbole \textquote{\art{}} für Artikel \textquote{\para{}} für Paragraphen, \absa{} für Absätze, \textquote{\satz} für Sätze, \textquote{\nr} für Nummern, \textquote{\lit} für Buchstaben.

\textquote{Die Würde des Menschen ist unantastbar} wird somit als \art{1} \absa{1} \satz{1} \GG zitiert.

\ankinote{zitieren}{Zitieren}{Grob-Genau, Gesetz; Art, §, Abs., S. Nr. lit.}
\end{document}